\scanpage{84}
% L1
\[
\unsure{p|\wp}(u+v)+pu+pv=\frac14\left(\frac{p'u-p'v}{pu-pv}\right)^2
\tag*{(21)}
\]

% L2–3
(en dérivant en $u$ et $v$ l'expression de $\frac{p'u}{pu-pv}$, et y remplaçant $p''$ en fonction de $p$)

% L4
faisant $v=u$

% L5
\[
p(2u)+2pu=\frac14\left(\frac{p''u}{p'u}\right)^2
\tag*{(21 bis)}
\]

% L6
L'application de la formule d'Hermite à $p(2u)$ donne

% L7
\[
4p(2u)=pu+p(u+\omega_1)+p(u+\omega_2)+p(u+\omega_3)
\tag*{(22)}
\]

% L8, boxed heading
5) \fbox{Fonctions $pu-e_i$}

% L9–10
L'application de la formule d'Hermite à $pu-e_i$ donne

% L11–12
\[
\frac{1}{pu-e_i}=\frac{p(u+\omega_i)-e_i}{(e_j-e_i)(e_k-e_i)}
\quad(i,j,k\ \text{\unsure{permutation} de }1,2,3)
\tag*{\textcircled{\scriptsize 23}}
\]

% L13
(compte tenu de $e_1+e_2+e_3=0$)

% L14–15
La racine carrée de \deleted{\unsure{cette|cet}} fonction est analytique méromorphe; la formule d'addition $pu-pv\inserted{\unsure{\ldots|=\ldots}}$ donne

% L16
\[
pu-e_i=-\frac{\sigma(u+\omega_i)\,\sigma(u-\omega_i)}{\sigma^2u\,\sigma^2\omega_i}
=\left(\frac{\sigma_i u}{\sigma u}\right)^{\unsure{2|7}}\quad\text{où}
\tag*{\textcircled{\scriptsize 24}}
\]

% L17
\[
\sigma_i u=\frac{1}{\sigma\omega_i}\,\sigma(u+\omega_i)e^{-\eta_i u}
\tag*{\textcircled{\scriptsize 25}}
\]

% L18–21, foot
On trouve que $\sqrt{pu-e_i}=\frac{\sigma_i u}{\sigma u}$ a les périodes $2\omega_i$ et $4\omega_j$ (mais non $2\omega_i$ et $2\omega_j$) -- $i\ne j$; elle change de signe si on ajoute $2\omega_i$ à $u$. En d'autres termes

% lower left
% archive sheet number: 83
